\section{Formal model of defense leakage} \label{formalmodel}
We adopt a formal language for specifying the security guarantees of a \anna{comprehensive??} Spectre defense.
% , that allows us to generate guarantee-violating execution graphs and concretize them into litmus tests to do targeted testing on an implementation. 
We take inspiration from prior work to 
We survey the literature of hardware defenses to Spectre attacks and conclude that a formal Spectre defense specification should consist of three components: an execution contract (what is speculative), a protection set (what is protected and how protection propagates), and a leakage contract (which operands can transmit).
\begin{enumerate}
    \item \textbf{Leakage contract:}
A leakage contract is the set of pairs of \textit{transmitters} (instructions that leak one or more operands via a hardware side channel) and their unsafe operand(s), as recognized by the defense.
\item \textbf{Execution contract:} 
an execution contract characterizes what speculative control- and data-flow is allowed in the $\mu$arch by identifying which \textit{speculation primitives} (instructions that initiate speculative executions, e.g., a conditional branch) the defense considers. 
\item \textbf{Protection set:}
A protection set is the set of architectural state (memory and registers) that the defense promises to not leak architecturally and/or speculatively. Along with the protection set we encode a protection set update rule, that represents the updates that a given protection set undergoes with each instruction.
\end{enumerate}
We then construct a per-defense \textit{leakage predicate}.
A defense's \emph{leakage predicate} combines its three components: an execution violates it when data in the protection set reaches, through rf and ddi edges, an operand in the leakage contract of an instruction that is speculative under the execution contract.
% \textbf{1. Leakage contract}
% A leakage contract is the set of pairs of \textit{transmitters} (instructions that leak one or more operands via a hardware side channel) and their unsafe operand(s), as recognized by the defense.
% \textbf{2. Execution contract}
% An execution contract characterizes what speculative control- and data-flow is allowed on the microarchitecture by identifying what \textit{speculation primitives} (instructions that initiate speculative executions, e.g., a conditional branch) the defense considers. 
% \textbf{3. Protection set}
% A protection set is the set of architectural state (memory and registers) that the defense promises to not leak architecturally and/or speculatively. Along with the protection set we encode a protection set update rule, that represents the updates that a given protection set undergoes with each instruction.
% todo these components sufficient for now etc etc 

\subsection{Execution Syntax}
\input{tikz/execvocab_v2_8pt}

\auditnote{fact}{fact}{\textbf{About the caption of Table~\ref{tab:execvocab}:} The vocabulary is shared, but more than the defense block differs. Recon\_6 and SPT\_6\_oneLP lack rf\_from\_most\_recent\_writer. SPT\_6\_oneLP rewrites skeleton code: operand-typed a and last\_committed\_protset\_p, secure\_speculation\_scheme\_p, tag\_xm, and an exactly-one-leak-point run predicate. Recon\_6 changes a skeleton call site to hardware\_protection\_policy[p]. This is open as a DECIDE in URGENT\_TODO.md:31.  \textbf{PROPOSED TEXT} (edits the sentence just before this note): The defense block (speculation contract, protection set, leakage function) is what defines each defense; the models also differ in a few skeleton facts (Section\~{}\textbackslash{}ref\{...\}).  \textit{Source:} diff STT\_6.als vs Recon\_6.als (rf fact; :229) and vs SPT\_6\_oneLP.als (:220-284, 313-314); URGENT\_TODO.md:31}
\notebr
\auditnote{fact}{fact}{\textbf{Update to the note above (2026-11-09):} two of its differences are gone. Every in-use model now has rf\_from\_most\_recent\_writer, and the production SPT model (SPT\_6\_oneLP\_v2, approved 2026-11-08) asks for an ordinary violation (at least one leaking operand). Still different outside the defense block: SPT tracks its protection set per operand, not per state (a, last\_committed\_protset\_p, secure\_speculation\_scheme\_p, tag\_xm), and Recon\_6 passes the perturbation into hardware\_protection\_policy. The ultimate models (STT\_6\_u approved for the clean runs on 2026-11-09) put every defense block on one shared skeleton with per-operand protection; once they are the evaluated models, the shared-skeleton claim is exact. \newline \textbf{PROPOSED TEXT} (edits the sentence before the note above): The defense block (speculation contract, protection set, leakage function) defines each defense; the rest of each model is shared, except that SPT tracks its protection set per operand. \newline \textit{Source:} STAGE1\_alloy/models/Recon\_6.als:163, :239; SPT\_6\_oneLP\_v2.als:6-17, :184, :254-256, :346; STAGE1\_alloy/models/proposed/ultimate/SKELETON\_u.als; URGENT\_TODO.md decision U1}
\notebr
To express a leakage predicate distinguishing leaky from non-leaky executions (e.g., the load-transmitter leakage in Figure~\ref{fig:ddi-rf}), we model program executions as directed happens-before graphs, inspired by MCM analyses~\cite{herdingcats,lustig}\anna{add citations}, and define leakage as a predicate over them, as in \ref{formalmodel}.
% \figtodo{cite the basic execution-graph example, Fig.~\ref{fig:ddi-rf} (the dataflow figure, redrawn in the new notation), here; the violation definition below says I3, the figure now says i3}
Generated tests consist of a total program order
\urgentbr
\urgentnote{\textbf{VOCAB}: \textbf{ADD DEFINITION} (new text, insert here): \emph{spo}, the speculative program order, the total order on instructions.}
\notebr
of \textbf{Instruction} nodes, typed as transmitters (Load, Store, Branchx, Otherx) or non-transmitters (Branchn, Othern). Input/output \textbf{Operands} (inmem, inreg, inaddr, outreg, outmem) refer to \textbf{State} objects (Reg, Mem). Instructions carry $\mu$arch flags for committed/resolved status
\urgentbr
\urgentnote{\textbf{VOCAB}: \textbf{ADD DEFINITION} (new text, insert here): \emph{resolved} (a branch's outcome, or with store bypass a store's address, is known) and \emph{committed} (retired; implies resolved; a prefix of spo).}%
\urgentbr
\urgentnote{\textbf{VOCAB}, replaces the definition just above: \emph{resolved} (the defense knows a branch's outcome, or with store bypass a store's address; a branch that has executed but not yet squashed its wrong path still counts as unresolved, as decided 2026-11-08) and \emph{committed} (retired; implies resolved; a prefix of spo).}%
\notebr
, enforced in simulation to replicate execution.
\notebr
\auditnote{fact}{fact, low}{Only resolution is actively enforced, by the branch force, resolve stall and STL hold. Commit status is measured, not enforced. The evaluated fixed configs set the fnc commit stall to 0, and the checkers only require lc to retire and the signal to come before fnc retires. The Gem5 section says this correctly ('passively measuring'). \newline \textbf{PROPOSED TEXT} (edits the sentence just before this note): Instructions carry \$\textbackslash{}mu\$arch flags for committed/resolved status; resolution is enforced in simulation and commit status is checked against the trace. \newline \textit{Source:} runconfigs/fixed\_20261030/*.jsonc (fnc\_commit\_stall\_cycles 0); STAGE3\_gem5/check\_ld.py:95-114}
\notebr
Data-dependence edges (rf, ddi)
\urgentbr
\urgentnote{\textbf{VOCAB}: \textbf{ADD DEFINITION} (new text, insert here): \emph{rf} (reads-from: an output operand to an input operand of a later instruction on the same location) and \emph{ddi} (data-dependence internal: an input to an output of the same instruction).}
\notebr
track information flow to transmitters. First-order axioms constrain legal edge placements. We choose the instructions, operands, and state that we use because we find it to be the minimal set of state needed to represent the distinctions in security guarantees in the body of Spectre defenses we study. Figure \ref{tab:execvocab} contains the execution graph syntax.
\notebr
\auditnote{fig}{figure, low}{This prints 'Figure I'; the target is a table. \newline \textbf{PROPOSED TEXT} (edits the sentence just before this note): Table\~{}\textbackslash{}ref\{tab:execvocab\} contains the execution-graph syntax. \newline \textit{Source:} p.4 left column, last line of the Execution Syntax paragraph.}
\notebr
% \figtodo{new vocabulary table, Table~\ref{tab:execvocab}: reference it here; the prose above can shrink to one sentence (suggested text in simspect paper\_figures/FIGURES.md)}

A \textit{violation} is defined as
\notebr
\textit{a data-dependence path from a protection-set state element as updated by the most recently committed instruction,
\auditnote{fact}{fact, high}{No model updates the protection set only up to the last committed instruction. last\_committed\_protset\_p applies the update rule to every instruction in spo order, committed or not. STT's rule removes the memory read by any load with no older unresolved branch. The commit-based rule is commented out. The dataflow figure caption already states the resolution-based rule, so the text and that caption disagree.  \textbf{PROPOSED TEXT} (edits the sentence just before this note): a data-dependence path from a protection-set state element (after the defense's update rule has been applied along speculative program order) to a leaked operand of a transmitter, ...  \textit{Source:} STAGE1\_alloy/models/STT\_6.als:232-243 (fold over f..f.spo\^{}6), :329-331 (rule), :330 (commit rule commented out); tikz/dataflow\_8pt.tex caption}
\auditnote{fact}{fact, high}{\textbf{Update (2026-11-09):} the ultimate STT model, STT\_6\_u (STT\_6\_v2 semantics), was approved for the clean runs. It taints the value read by every unsafe load, forwards taint along rf and ddi, and untaints only by the speculation contract, with a per-operand protection set. If the paper's results come from it, define a violation as a data-dependence path from a value read by a speculative load to a leaked operand of a speculative transmitter, and drop ``protection-set state element''. \textit{Source:} URGENT\_TODO.md decision U1 and the STT\_6\_u DECIDE; STAGE1\_alloy/models/proposed/ultimate/STT\_6\_u.als}
to a leaked operand of a transmitter, where the transmitter is in the leakage contract and speculative as per the execution contract.} In Figure~\ref{fig:ddi-rf}, the transmitter load in \texttt{I3} has an address dependency on the load in \texttt{I1}
\notebr
\auditnote{fig}{figure, low}{The text names the instructions I1, I2, I3 in typewriter, but Fig. 2 now labels them i1, i2, i3. \newline \textbf{PROPOSED TEXT} (edits the sentence just before this note): ... the transmitter load i3 has an address dependency on the load i1 through the ALU op i2 ... \newline \textit{Source:} p.4 text vs Fig. 2 row labels.}
\notebr
through \texttt{I2} arithmetic, and the memory is in the protection set, creating a violation. 

\input{tikz/dataflow_8pt}% redrawn basic example; old tikz/alloyviolation kept, 2026-11-02
\input{tikz/pipeline_full_8pt1col}
\notebr
\auditnote{fig}{figure}{\textbf{About the caption of Fig.~\ref{fig:pipeline}:} 'Non-invasive' (also in the caption) overstates the harness: the directives change prediction and timing on purpose (BTB force, resolve stall, fnc stall). The 'selects hooks' arrow from transmission mechanisms into the harness has no code behind it.  \textbf{PROPOSED TEXT} (edits the sentence just before this note): "+ Minimal harness (defense logic untouched): probes; speculation and timing directives"; drop the 'selects hooks' arrow, or label it 'design time'.  \textit{Source:} BIGFIG\_FACT\_AUDIT.md rows 30 and 39, decisions D6 and D8; pipeline.tex gem5-hooks paragraph describes the stalls as active control hooks.}
\notebr
\auditnote{fig}{figure}{\textbf{Update to the note above (2026-11-09):} the `selects hooks' arrow stays (decided 2026-11-02), but the caption or text must say the hooks are chosen at design time: no code passes the transmission mechanism to gem5 at run time. `Non-invasive' still overstates the harness: the branch force, the resolve stall and the store-address hold change prediction and timing on purpose, and the CleanupSpec and SpecLFB runs start each test with a hot or cold cache at random (decided 2026-11-09); only the defense logic is untouched (wording still open). \newline \textbf{PROPOSED TEXT}: ``+ Minimal harness (defense logic untouched): probes; speculation and timing directives''; in the caption, ``\ldots{} under a minimal harness whose hooks are chosen at design time from the defense's transmission mechanisms''. \newline \textit{Source:} URGENT\_TODO.md (pipeline figure decisions; cache-setup decision 2026-11-09); wiki/design/RESOLUTION\_HOOKS.md sec.~1; runconfigs/twostage\_fv2\_20261109/E13\_STT\_6\_fv2.jsonc, E14 (cache\_setup random)}
\notebr
\auditnote{fig}{figure, high}{\textbf{About the caption of Fig.~\ref{fig:pipeline}:} No evaluated enumeration reached UNSAT; every testset stops at max\_instances (100k), and Sec. V's OTB paragraph says the shape is 'not reached within its instance cap'. Fig. 3 calling the output 'bounded-exhaustive' contradicts the paper's own results.  \textbf{PROPOSED TEXT} (edits the sentence just before this note): "Minimal leaking execution graphs, up to the bound and instance cap" (or keep 'bounded-exhaustive' only if a run reaches UNSAT; decision D2).  \textit{Source:} BIGFIG\_FACT\_AUDIT.md rows 7, 9 and D2; STAGE1\_alloy/models/MANIFEST.md (100k cap); results.tex sec:bug-otb 'does not reach this shape within its instance cap'.}
\notebr
\auditnote{fig}{figure, high}{\textbf{Update to the note above (2026-11-09):} `Bounded-exhaustive' now holds in part. In the two-stage runs, Stage~1a (up to six instructions) reaches UNSAT for STT and ReCon and their store-bypass variants, but SPT's Stage~1a does not converge: its four-instruction partition is capped at 200,000 graphs. The June numbers come from flat test sets that stop at their instance cap. \newline \textbf{PROPOSED TEXT}: ``Bounded-exhaustive (STT, ReCon) set of minimal violating execution graphs'', or keep the label once SPT's Stage~1a is settled. \newline \textit{Source:} wiki/state/EVAL\_MANIFEST.md `Stage 1' section}
\notebr
\auditnote{fig}{figure, high}{\textbf{About the caption of Fig.~\ref{fig:pipeline}:} The caption calls Stage 1b's output minimal, but the figure itself says Stage 1b loosens minimality (slab 2'). Also, per the fact audit, Stage 1b is an STT-only prototype that produced none of the reported results, and neither the figure nor the caption says so.  \textbf{PROPOSED TEXT} (edits the sentence just before this note): Stage\~{}1 uses Alloy to enumerate minimal leaking execution graphs up to a bound; Stage\~{}1.5 (prototype, not used for the results below) re-enumerates with loosened minimality, pinned to each Stage\~{}1 graph.  \textit{Source:} Fig. 3 slab '2'. Loosened minimality requirements'; BIGFIG\_FACT\_AUDIT.md row 10 and decision D1 (no campaign used Stage 1b).}
\notebr
\auditnote{fig}{figure, high}{\textbf{Update to the note above (2026-11-09):} Stage~1b is no longer a prototype. It runs one pinned Alloy run per Stage-1a graph, in parallel, at exactly six instructions, keeps RR, RC and RS but not RI, and caps each graph's flesh-out (about 300,000 graphs per defense); the clean-run results will come from it, no number in the paper does yet. The caption still calls its output minimal. \newline \textbf{PROPOSED TEXT}: Stage~1a uses Alloy to enumerate minimal violating execution graphs of up to six instructions; Stage~1b grows each one into six-instruction graphs under loosened minimality, a capped number per Stage-1a graph. \newline \textit{Source:} wiki/state/EVAL\_MANIFEST.md `Stage 1.5' section; STAGE1\_alloy/models/proposed\_twophase/STT\_6\_core\_ri.als:289-310}
\notebr
\auditnote{fact}{fact}{\textbf{About the caption of Fig.~\ref{fig:pipeline}:} The harness changes prediction and timing on purpose. It forces branch directions, holds branch resolution and store addresses, and on the resolve-hook builds wakes the CPU at release. Only the defense logic is untouched. The same word appears in the figure body and in pipeline.tex ('non-invasively instrument'). The wording is an open item in URGENT\_TODO.md:44.  \textbf{PROPOSED TEXT} (edits the sentence just before this note): Stage\~{}3 runs them on the defense's gem5 build under a minimal harness that leaves the defense logic untouched  \textit{Source:} wiki/design/RESOLUTION\_HOOKS.md sec.1 (behaviour changes; edit 4); gem5-hooks/README.md (STL hooks); paper\_figures/plans/BIGFIG\_FACT\_AUDIT.md row 30; URGENT\_TODO.md:44}
\notebr
\auditnote{fig}{figure, low}{\textbf{About the caption of Fig.~\ref{fig:pipeline}:} The naming is inconsistent: Fig. 3 says 'Inspector' and 'Pipeline info artifact', while the text and the Fig. 7 caption say 'checker', 'per-transmitter pipeline predicate' and 'gem5 execution trace'. Stage 1a is 'Abstract test generation' and Stage 1b 'Abstract test enumeration', which implies a difference that does not exist.  \textbf{PROPOSED TEXT} (edits the sentence just before this note): Use one name for the checker across Fig. 3, Fig. 7 and the text, and 'Pipeline trace' for the artifact; e.g. Stage 1a 'Minimal-graph enumeration', Stage 1b 'Graph expansion'.  \textit{Source:} Fig. 3 vs Fig. 7 caption ('How the checker decides') and pipeline.tex Transmission Requirement; BIGFIG\_FACT\_AUDIT.md rows 1, 35.}
\notebr
\auditnote{fig}{figure}{\textbf{About the caption of Fig.~\ref{fig:pipeline}:} The caption explains only blue. It never explains the tan, lavender and light-red slabs 3a-3c (which carry the same meaning as in the graph figures), the dashed 'Same as Stage 1a' ghost slabs, the + and prime marks, the numbered chevrons, or the dashed 'SimSpect goal' arrow.  \textbf{PROPOSED TEXT} (edits the sentence just before this note): Blue marks defense-specific components; slabs 3a-3c use the speculation (tan), protection-set (lavender) and transmitter (red) colors of the graph figures; dashed slabs are unchanged from Stage 1a, + marks an added and ' a changed constraint.  \textit{Source:} p.4 Fig. 3 render; VISUAL\_LANGUAGE.md CURRENT RULES (pipeline 3a tan, 3b lavender, 3c light red).}
\notebr

\subsection{Example contract}
\notebr
\urgentbr\urgentnote{\textbf{RENAME (Anna, 2026-11-09): call the two enumeration stages of Fig.~\ref{fig:pipeline} Stage~1a and Stage~1b, not Stage~1 and Stage~1.5.} In tikz/pipeline\_full\_8pt1col.tex: the slab titles (``\textbackslash{}textbf\{Stage 1\}'' $\to$ ``Stage 1a'', ``Stage 1.5'' $\to$ ``Stage 1b''), the two ``Same as Stage 1'' ghost labels (``Same as Stage 1a''), and the caption sentence ``Stages~1 and~1.5 use Alloy'' $\to$ ``Stages~1a and~1b''. Stage~2 and Stage~3 keep their numbers. The same renaming applies wherever the text says Stage~1 or Stage~1.5 for these two: results.tex (the Coverage item, ``Stage-1 graphs'' / ``Stage-1.5 graphs''), pipeline.tex, intro.tex, discussion.tex, and the other notes in this file. The other pipeline variants (tikz/pipeline\_full.tex, pipeline\_full\_8pt.tex) are unused; rename there only if you bring them back. \textbf{Status (2026-11-11):} the figure and its caption are renamed, every note now says Stage~1a/1b, and STAGE NAMES notes in pipeline.tex mark the prose to change (the stage walk-through near the top of the section, the Test Generation subsubsection, the interleave caption). The hidden to-do lists in results.tex and discussion.tex carry a comment. Once those are done, delete this note.}
% scope of the types of defenses we can represent
For example, we define the defense specification for Speculative Taint Tracking, or STT, a defense that prevents leakage by delaying 'tainted' transmitters. 
\begin{enumerate}
    \item \textbf{Leakage contract:}
\verb|(load, load address)|, \verb|(branch, condition)|
\item \textbf{Execution contract:} 
An instruction is speculative when it is uncommitted and there are older, unresolved branches. 
\item \textbf{Protection set:}
All memory; a load with no older unresolved branch removes the location it reads from the protection set.
\end{enumerate}
% STT's leakage function is (load, load address), (branch, input register), and (reg-op transmitter, input address). STT does not consider stores to transmit any operands. STT's execution contract can be set as either that an instruction is non-speculative upon commit, or that an instruction is non-speculative when there exists no older unresolved branches. STT's protection set is all memory, with no protection set updates. 

% \subsection{Execution contracts in practice}
% \subsection{Scope of defenses}
% \begin{wrapfigure}{r}{.5\textwidth}
%   \begin{center}
%     \includegraphics[width=.5\textwidth]{images/examplefigure.png}
%   \end{center}
%   \caption{We generate a test from STT using our pipeline.}

%   \label{fig:pipeline}
  
% \end{wrapfigure}
% \begin{figure}[h]
%   \centering
%   \includegraphics[width=0.4\textwidth]{images/examplefigure.png}
%   \caption{We generate a test from STT using our pipeline.}
%   \label{fig:pipeline}
% \end{figure}